\documentclass[a4paper,11pt]{article}

% ════════════════════ Kodowanie i język ════════════════════
\usepackage[T1]{fontenc}
\usepackage[utf8]{inputenc}
\usepackage{lmodern}
\usepackage[polish]{babel}

% ════════════════════ Marginesy ════════════════════
\usepackage[margin=2.5cm]{geometry}

% ════════════════════ Matematyka ════════════════════
\usepackage{amsmath}
\usepackage{amssymb}
\usepackage{siunitx}
\sisetup{
	output-decimal-marker = {,},
	per-mode = symbol,
	range-phrase = { -- },
}
\DeclareSIUnit{\baud}{Bd}
\DeclareSIUnit{\pkt}{pkt}
\DeclareSIUnit{\PLN}{PLN}

% ════════════════════ Grafika i tabele ════════════════════
\usepackage{graphicx}
\usepackage{float}
\usepackage{booktabs}
\usepackage{array}
\usepackage{tabularx}
\usepackage{longtable}
\usepackage{caption}

% ════════════════════ Kolory i linki ════════════════════
\usepackage[dvipsnames]{xcolor}
\usepackage{hyperref}
\hypersetup{
	colorlinks=true,
	linkcolor=NavyBlue,
	citecolor=NavyBlue,
	urlcolor=RoyalBlue,
}

% ════════════════════ Listingi kodu ════════════════════
\usepackage{listings}
\lstset{
	language=C,
	basicstyle=\ttfamily\footnotesize,
	keywordstyle=\color{NavyBlue}\bfseries,
	commentstyle=\color{Gray}\itshape,
	stringstyle=\color{OliveGreen},
	numbers=left,
	numberstyle=\tiny\color{Gray},
	frame=single,
	breaklines=true,
	captionpos=b,
	tabsize=4,
	showstringspaces=false,
}

% ════════════════════ Nagłówek / stopka ════════════════════
\usepackage{fancyhdr}
\pagestyle{fancy}
\fancyhf{}
\fancyhead[L]{\small Bezkontaktowy miernik odległości --- Etap 2 (prototyp)}
\fancyhead[R]{\small Wygraj Indeks WETI 2026}
\fancyfoot[C]{\thepage}
\renewcommand{\headrulewidth}{0.4pt}

\usepackage{setspace}
\onehalfspacing
\usepackage{parskip}
\usepackage{enumitem}

% ══════════════════════════════════════════════════════════════
\begin{document}

% ════════════════════ STRONA TYTUŁOWA ════════════════════
\begin{titlepage}
	\centering
	\vspace*{1.5cm}

	{\Large Ogólnopolski Konkurs Projektowy}\\[0.3cm]
	{\Large\textbf{,,Wygraj Indeks WETI'' --- edycja 2026}}\\[0.2cm]
	{\large Wydział Elektroniki, Telekomunikacji i~Informatyki\\
		Politechnika Gdańska}

	\vspace{1.5cm}

	{\Huge\bfseries Bezkontaktowy miernik odległości}\\[0.5cm]
	{\Large Metoda hybrydowa Pulse-Phase Ranging\\z~kompensacją temperatury}\\[0.4cm]
	{\Large \textbf{Etap 2 --- dokumentacja prototypu}}\\[0.3cm]
	{\large Realizacja toru zgrubnego (Time-of-Flight) ze~sprzętową,\\
		bezlatencyjną detekcją echa}

	\vspace{2.0cm}

	{\large
		\begin{tabular}{rl}
			\textbf{Autor:}    & Kacper Popko \\
			\textbf{Opiekun:}  & Mateusz Karanowski \\
			\textbf{Szkoła:}   & II Liceum Ogólnokształcące im.\ ks.\ Anny z~Sapiehów \\
			                   & Jabłonowskiej w~Białymstoku \\
			\textbf{E-mail:}   & kacper@kajpa.pl \\
			\textbf{Data:}     & Czerwiec 2026 \\
		\end{tabular}
	}

	\vspace{1.5cm}

	\begin{tabular}{ll}
		\toprule
		\textbf{Mikrokontroler} & ATtiny1616-SN (tinyAVR 1-series) \\
		\textbf{Metoda} & Pulse-Phase Ranging \SI{40}{\kilo\hertz}: ToF (obwiednia ADC) + faza (Goertzel) \\
		\textbf{Zakres} & \SIrange{10}{50}{\centi\metre} \\
		\textbf{Toolchain} & AVR-GCC 15.2 + avrdude 8.1 (czysto rejestrowo) \\
		\bottomrule
	\end{tabular}

	\vfill
\end{titlepage}

\tableofcontents
\newpage

% ══════════════════════════════════════════════════════════════
\section{Wprowadzenie i~zakres Etapu 2}
\label{sec:wstep}

Niniejszy dokument stanowi dokumentację \textbf{zbudowanego i~uruchomionego
prototypu} bezkontaktowego miernika odległości, dostarczanego w~ramach Etapu~2
konkursu ,,Wygraj Indeks WETI 2026''. W~odróżnieniu od dokumentacji Etapu~1
(projekt koncepcyjny) opisuje on \emph{rzeczywisty stan urządzenia} po montażu,
uruchomieniu i~strojeniu na stanowisku pomiarowym.

Projekt kontynuuje przyjętą w~Etapie~1 \textbf{metodę hybrydową Pulse-Phase
Ranging z~kompensacją temperatury}. Metoda ta składa się z~dwóch etapów: pomiaru
\emph{zgrubnego} czasu przelotu (Time-of-Flight), który eliminuje
niejednoznaczność, oraz pomiaru \emph{dokładnego} fazy echa (algorytm
Goertzela), dającego rozdzielczość sub-milimetrową. W~prototypie
\textbf{zaimplementowano i~uruchomiono oba etapy}: pomiar ToF przez
próbkowanie obwiedni echa przetwornikiem ADC oraz wyznaczanie fazy nośnej
\SI{40}{\kilo\hertz} algorytmem Goertzela, a~także pełną kompensację
temperatury. Zapewnia to poprawny pomiar w~całym zakresie
\SIrange{10}{50}{\centi\metre} (sek.~\ref{sec:firmware}--\ref{sec:wyniki}).
Dodatkowo zweryfikowano sprzętową, bezlatencyjną ścieżkę detekcji progowej
\texttt{AC1+DAC1+EVSYS+TCB0} (sek.~\ref{sec:hw_capture}).

W~trakcie realizacji prototypu wprowadzono szereg zmian względem założeń
projektowych --- wynikających z~praktycznej weryfikacji sprzętu oraz z~analizy
dokumentacji mikrokontrolera. Zmiany te zestawiono w~sekcji~\ref{sec:zmiany}.
Najistotniejsza dotyczy ścieżki detekcji echa: pierwotnie zakładano użycie
komparatora \texttt{AC0}, jednak weryfikacja Tabeli~5-1 noty katalogowej
wykazała, że wejście echa (pin \texttt{PB4}) jest dodatnim wejściem
\texttt{AINP3} komparatora \texttt{AC1}, a~nie \texttt{AC0}. Skorygowano to
w~firmware, dzięki czemu sprzętowa, bezlatencyjna detekcja czasu przelotu
działa poprawnie.

\subsection{Spełnienie wymagań konkursowych}

\begin{table}[H]
	\centering
	\caption{Status wymagań konkursowych w~prototypie}
	\label{tab:wymagania}
	\begin{tabularx}{\textwidth}{llX}
		\toprule
		\textbf{Parametr} & \textbf{Wymóg} & \textbf{Stan w~prototypie} \\
		\midrule
		Napięcie zasilania & \SI{5}{\volt} & Spełnione; LDO MCP1700 $\to$ \SI{3,3}{\volt} \\
		Zakres pomiaru & \SIrange{10}{50}{\centi\metre} & Spełnione; potwierdzone pomiarami \\
		Wykrywany obiekt & Karton $20\times20$~\si{\centi\metre} & Spełnione \\
		Błąd RMSE & im mniejszy tym lepiej & $\approx\SI{6}{\milli\metre}$ (metoda ToF, sek.~\ref{sec:wyniki}) \\
		Pobór mocy & $100/P_{[\si{\milli\watt}]}$ & Standby+RTC, $P_{avg}<\SI{1}{\milli\watt}$ (sek.~\ref{sec:moc}) \\
		Koszt elementów & $2500/C_{[\text{PLN}]}$ & \SI{43,73}{\PLN} $\to$ 50~pkt (bez~zmian) \\
		\bottomrule
	\end{tabularx}
\end{table}

% ══════════════════════════════════════════════════════════════
\section{Architektura zbudowanego urządzenia}
\label{sec:arch}

\subsection{Schemat blokowy toru sygnałowego}

Tor pomiarowy realizuje pomiar czasu przelotu (Time-of-Flight) impulsu
ultradźwiękowego \SI{40}{\kilo\hertz}:

\begin{enumerate}[nosep]
	\item Mikrokontroler generuje paczkę (burst) \SI{40}{\kilo\hertz} na pinie
	\texttt{PB0}, wzmacnianą tranzystorem \texttt{BC817} sterującym nadajnik US.
	\item Echo odbierane przez przetwornik RX jest wzmacniane dwustopniowym
	torem na \texttt{MCP6002} (wzmocnienie $\sim$10$\times$, filtr pasmowy wokół
	\SI{40}{\kilo\hertz}).
	\item Wzmocniony sygnał echa (pin \texttt{PB4 = ADC0 AIN9}) jest szybko
	próbkowany przetwornikiem ADC; z~obwiedni wyznaczany jest początek echa
	(czas przelotu), a~z~przebiegu --- faza nośnej (Goertzel).
	\item Niezależnie ten sam sygnał trafia na komparator \texttt{AC1}
	(próg z~\texttt{DAC1}) $\to$ \texttt{EVSYS} $\to$ \texttt{TCB0} ---
	zweryfikowana, sprzętowa, bezlatencyjna ścieżka detekcji progowej.
	\item Temperatura z~termistora NTC koryguje prędkość dźwięku;
	odległość $D = c(T)\cdot t_{ToF}/2$.
\end{enumerate}

\subsection{Mapa pinów ATtiny1616 (stan rzeczywisty)}
\label{sec:pinout}

Tabela~\ref{tab:pinout} przedstawia \textbf{faktyczne} przyporządkowanie pinów
w~zbudowanym prototypie. Różni się ono od wstępnego projektu (Etap~1) ---
zmiany wynikły z~optymalizacji rozprowadzenia ścieżek PCB oraz z~korekty
przypisania komparatora (patrz sek.~\ref{sec:zmiany}).

\begin{table}[H]
	\centering
	\caption{Mapa pinów ATtiny1616 w~prototypie}
	\label{tab:pinout}
	\begin{tabularx}{\textwidth}{llX}
		\toprule
		\textbf{Pin} & \textbf{Funkcja} & \textbf{Opis / peryferiał} \\
		\midrule
		PA0 & UPDI & Programowanie / debug \\
		PA3 & PRZYCISK & Wejście, pull-up, aktywny stan niski (start pomiaru) \\
		PA4 & LED & Wyjście, sygnalizacja pomiaru \\
		PA5 & TERMISTOR & ADC0 AIN5 (kompensacja temperatury, NTC \SI{10}{\kilo\ohm}) \\
		PA6 & GATE Q1 & Sterowanie P-MOSFET BSS84 (zasilanie toru RX) \\
		PB0 & TX BURST & Generacja \SI{40}{\kilo\hertz} (programowa) $\to$ R1 $\to$ BC817 \\
		PB2 & USART0 TXD & Transmisja wyniku, \SI{115200}{\baud} 8N1 \\
		PB4 & ECHO & AC1 AINP3 (wyjście toru RX $\to$ komparator) \\
		\bottomrule
	\end{tabularx}
\end{table}

% ══════════════════════════════════════════════════════════════
\section{Oprogramowanie mikrokontrolera}
\label{sec:firmware}

Firmware napisano w~języku C, kompilowany \textbf{AVR-GCC 15.2} (bez frameworka
Arduino), wgrywany przez \textbf{avrdude 8.1} interfejsem UPDI. Konfiguracja
peryferiów odbywa się na poziomie rejestrowym. Plik źródłowy:
\texttt{firmware/main.c}; plik wynikowy: \texttt{firmware/main.hex}
(${\sim}$\,\SI{6,4}{\kilo\byte} Flash z~\SI{16}{\kilo\byte}, \SI{332}{\byte} RAM).

\subsection{Pomiar ToF przez próbkowanie obwiedni (ADC)}
\label{sec:adc}

Właściwy pomiar odległości realizowany jest przez \textbf{szybkie próbkowanie
sygnału echa przetwornikiem ADC} (pin \texttt{PB4 = ADC0 AIN9}, tryb 8-bit,
prescaler DIV4, ${\sim}\SI{8,7}{\micro\second}/\text{próbkę} \approx
\SI{115}{ksps}$, tj.\ ${\sim}2{,}9$ próbki na okres \SI{40}{\kilo\hertz}).
Minimalna pętla próbkująca zapisuje \SI{320}{} próbek do bufora; czas trwania
okna jest mierzony licznikiem \texttt{TCB0} (\SI{100}{\nano\second}/takt), co
daje dokładną podstawę czasu.

W~post-processingu wyznaczany jest \textbf{początek echa (onset)} --- pierwsza
próbka przekraczająca próg \texttt{ONSET\_DEV} --- którego czas odpowiada
rzeczywistemu czasowi przelotu: $D = c(T)\cdot t_{onset}/2$. Metoda obwiedni
jest odporna na amplitudę (wykrywa też słabe echa z~\SI{40}{\centi\metre}),
a~zapisany bufor próbek służy następnie etapowi fazowemu (Goertzel).
Detekcja przez ADC okazała się czulsza od progowej ścieżki sprzętowej
(sek.~\ref{sec:hw_capture}) i~dodatkowo dostarcza przebiegu potrzebnego do
pomiaru fazy --- dlatego przyjęto ją jako podstawową.

\subsection{Pomiar fazy --- algorytm Goertzela}
\label{sec:goertzel}

Na buforze próbek echa uruchamiany jest \textbf{filtr Goertzela} dostrojony do
\SI{40}{\kilo\hertz}, liczący część rzeczywistą i~urojoną prążka DFT, a~stąd
amplitudę i~fazę nośnej. Faza odniesiona do \emph{startu burstu} (przez
dodanie członu $f_0\cdot t_{okna}$) koduje pozycję obiektu w~obrębie
$\lambda/2 \approx \SI{4,33}{\milli\metre}$:
\begin{equation}
	\frac{D}{\lambda/2} = f_0\,t_{ToF}
	= f_0\,t_{okna} + \frac{\varphi_0 - \varphi}{2\pi}.
\end{equation}
Wynik łączony jest ze~zgrubnym ToF: $D = n\cdot(\lambda/2) + D_{faza}$, gdzie
$n$ wynika z~pomiaru zgrubnego, z~kontrolą spójności (odrzucenie skoków o~pół
fali). Pomiarowo zweryfikowano, że faza zmienia się ${\sim}83^\circ/\si{\milli\metre}$
zgodnie z~teorią i~jest powtarzalna przy stałej odległości.

\subsection{Zweryfikowana sprzętowa ścieżka detekcji (AC1+TCB0)}
\label{sec:hw_capture}

Niezależnie zaimplementowano i~zweryfikowano \textbf{w~pełni sprzętową,
bezlatencyjną} ścieżkę detekcji progowej, nieobciążającą rdzenia CPU. Łańcuch
zdarzeń (datasheet sek.~14 i~29):

\begin{table}[H]
	\centering
	\caption{Konfiguracja rejestrów ścieżki detekcji echa}
	\label{tab:rejestry}
	\footnotesize
	\begin{tabularx}{\textwidth}{llX}
		\toprule
		\textbf{Peryferiał} & \textbf{Konfiguracja} & \textbf{Znaczenie} \\
		\midrule
		VREF & \texttt{CTRLC: DAC1REFSEL=2V5} & Referencja \SI{2,5}{\volt} dla DAC1/AC1 ($\leq V_{DD}$) \\
		     & \texttt{CTRLB: DAC1REFEN=1} & Wymuszone włączenie referencji \\
		DAC1 & \texttt{DATA=255, ENABLE=1} & Próg detekcji $\approx\SI{2,49}{\volt}$ (tuż pod $V_G$) \\
		AC1  & \texttt{MUXPOS=AINP3 (PB4)} & Wejście dodatnie = echo \\
		     & \texttt{MUXNEG=DAC} & Wejście ujemne = próg DAC1 \\
		     & \texttt{ENABLE=1, HYSMODE=10mV} & Histereza \SI{10}{\milli\volt} \\
		EVSYS & \texttt{ASYNCCH0=AC1\_OUT} & Źródło: wyjście AC1 \\
		      & \texttt{ASYNCUSER0=ASYNCCH0} & Odbiorca: TCB0 (input capture) \\
		TCB0  & \texttt{CNTMODE=CAPT, CAPTEI=1} & Zatrzask licznika na zboczu zdarzenia \\
		      & \texttt{CLKSEL=CLKDIV1} & \SI{10}{\mega\hertz}, \SI{100}{\nano\second}/takt \\
		\bottomrule
	\end{tabularx}
\end{table}

Próg detekcji jest ustawiony \emph{tuż poniżej} napięcia spoczynkowego toru
$V_G\approx\SI{2,69}{\volt}$. W~stanie spoczynku $V_{PB4}>V_{DAC1}$, więc
\texttt{AC1} ma stabilny stan wysoki. Echo wywołuje oscylacje, których
ujemne półfale schodzą poniżej progu --- każde przejście w~górę generuje
zbocze narastające, zatrzaskiwane przez \texttt{TCB0}. Pierwsze prawidłowe
zatrzaśnięcie po czasie blankingu odpowiada nadejściu echa.

\subsection{Algorytm pomiaru}

\begin{enumerate}[nosep]
	\item Wybudzenie z~uśpienia (Standby), włączenie toru analogowego (Q1 ON),
	stabilizacja \SI{5}{\milli\second}.
	\item Odczyt temperatury z~NTC (ADC 10-bit), obliczenie
	$c(T)=331{,}3+0{,}606\,T$ [m/s].
	\item Burst TX $N=10$ cykli \SI{40}{\kilo\hertz} (\SI{12,5}{\micro\second}
	półokres $\to$ dokładny rezonans przetwornika).
	\item Blanking (pominięcie ring-downu), szybkie próbkowanie obwiedni echa
	(ADC 8-bit, \SI{320}{} próbek) $\to$ onset $\to$ zgrubne $D$ (sek.~\ref{sec:adc}).
	\item Seria $M=9$ pomiarów, filtr medianowy $\to$ odporne $D_{coarse}$.
	\item Goertzel na buforze ostatniego pingu $\to$ faza $\to$ fuzja
	$D = n\cdot(\lambda/2)+D_{faza}$ z~kontrolą spójności (sek.~\ref{sec:goertzel}).
	\item Kalibracja liniowa $D = a\cdot D_{raw}+b$, transmisja przez USART,
	flush TX, powrót do uśpienia.
\end{enumerate}

\subsection{Kompensacja temperatury}

Termistor NTC \SI{10}{\kilo\ohm} ($B=3950$) w~dzielniku z~rezystorem
\SI{10}{\kilo\ohm} do masy. Rezystancja $R_{NTC}=R_{10}\cdot(1023-\text{ADC})/\text{ADC}$;
temperatura z~modelu Beta:
\begin{equation}
	\frac{1}{T} = \frac{1}{T_0} + \frac{1}{B}\ln\!\frac{R_{NTC}}{R_{25}},
	\qquad T_0 = \SI{298,15}{\kelvin}.
\end{equation}
Prędkość dźwięku skalowana liniowo z~temperaturą eliminuje dominujący błąd
systematyczny (zmiana $\sim\SI{0,17}{\percent}/\si{\kelvin}$).

% ══════════════════════════════════════════════════════════════
\section{Zarządzanie mocą}
\label{sec:moc}

Prototyp realizuje agresywną politykę uśpienia zapowiedzianą w~Etapie~1
(Scenariusz 2/3):

\begin{itemize}[nosep]
	\item \textbf{Tryb Standby} mikrokontrolera między pomiarami --- rdzeń CPU
	i~większość peryferiów zatrzymane.
	\item \textbf{Wybudzanie z~RTC/PIT} co \SI{125}{\milli\second} (licznik PIT
	taktowany wewnętrznym oscylatorem ULP \SI{32,768}{\kilo\hertz}); po wybudzeniu
	odpytywany jest przycisk.
	\item \textbf{Bramkowanie toru analogowego} tranzystorem Q1 (BSS84) ---
	wzmacniacz \texttt{MCP6002} jest zasilany wyłącznie podczas pomiaru
	($\sim\SI{30}{\milli\second}$), w~spoczynku pobiera \SI{0}{\micro\ampere}.
	\item \textbf{ADC i~AC} włączane na żądanie; nieużywane bufory wejściowe pinów
	wyłączone (\texttt{INPUT\_DISABLE}).
\end{itemize}

W~stanie spoczynku pobór prądu jest zdominowany przez oscylator ULP RTC
($\sim\SIrange{1}{3}{\micro\ampere}$). Przy typowym użyciu (pomiar na żądanie)
średni pobór mocy nie przekracza \SI{1}{\milli\watt}, co odpowiada maksymalnej
punktacji (\SI{50}{\pkt}) w~kryterium mocy.

\begin{table}[H]
	\centering
	\caption{Bilans prądowy (zmierzony / szacowany)}
	\label{tab:prad}
	\begin{tabularx}{\textwidth}{Xrr}
		\toprule
		\textbf{Stan} & \textbf{Prąd} & \textbf{Uwaga} \\
		\midrule
		Pomiar aktywny (op-amp + MCU + ADC) & $\sim\SI{4,3}{\milli\ampere}$ & zmierzone (\SI{5}{\volt}, INA219) \\
		Spoczynek bez uśpienia (tylko bramkowanie) & $\sim\SI{3,2}{\milli\ampere}$ & zmierzone \\
		Spoczynek Standby+RTC & $\sim\SIrange{1}{3}{\micro\ampere}$ & szac.\ (ULP osc + RTC) \\
		\bottomrule
	\end{tabularx}
\end{table}

% ══════════════════════════════════════════════════════════════
\section{Wyniki pomiarowe i~kalibracja}
\label{sec:wyniki}

Pomiary wykonano kartonem $20\times20$~\si{\centi\metre} ustawionym prostopadle
do osi przetworników, na czystej ścianie. Po kalibracji liniowej
($D=a\cdot D_{raw}+b$) uzyskano zgodność z~miarą wzorcową w~całym zakresie
\SIrange{10}{50}{\centi\metre}.

\begin{table}[H]
	\centering
	\caption{Przykładowa charakterystyka po kalibracji (metoda peak/ToF)}
	\label{tab:wyniki}
	\begin{tabular}{rrrr}
		\toprule
		\textbf{Odl.\ rzecz.\ [mm]} & \textbf{Odczyt śr.\ [mm]} & \textbf{Błąd [mm]} & \textbf{$\sigma$ [mm]} \\
		\midrule
		100 & 104 & $+4$  & 4{,}4 \\
		150 & 159 & $+9$  & 0{,}0 \\
		200 & 191 & $-9$  & 0{,}0 \\
		300 & 305 & $+5$  & 5{,}7 \\
		400 & 401 & $+1$  & 4{,}0 \\
		500 & 498 & $-2$  & 4{,}5 \\
		\bottomrule
	\end{tabular}
\end{table}

Błąd średniokwadratowy z~powyższych punktów:
\begin{equation}
	\text{RMSE} = \sqrt{\frac{1}{n}\sum_i (D_i - D_{i,ref})^2} \approx \SI{5,9}{\milli\metre}.
\end{equation}

Powtarzalność (odchylenie standardowe serii pomiarów dla stałej odległości)
wynosi $\sigma\approx\SI{4}{\milli\metre}$, a~liniowość charakterystyki
$R^2>0{,}999$. Podstawa czasu \texttt{TCB0} (\SI{100}{\nano\second}/takt)
nie ogranicza rozdzielczości --- limitują ją czynniki analogowe (czas
narastania obwiedni echa). Etap fazowy (Goertzel) zmienia się
${\sim}83^\circ/\si{\milli\metre}$ zgodnie z~teorią ($360^\circ$ na
$\lambda/2\approx\SI{4,33}{\milli\metre}$) i~jest powtarzalny przy stałej
odległości; pełne wykorzystanie do sub-mm ogranicza obecnie jitter pomiaru
zgrubnego (sek.~\ref{sec:dalsze}).

% ══════════════════════════════════════════════════════════════
\section{Zmiany względem Etapu 1}
\label{sec:zmiany}

\begin{table}[H]
	\centering
	\caption{Rozbieżności projekt (Etap 1) vs.\ prototyp (Etap 2)}
	\label{tab:zmiany}
	\begin{tabularx}{\textwidth}{p{3.2cm}XX}
		\toprule
		\textbf{Element} & \textbf{Etap 1 (projekt)} & \textbf{Etap 2 (prototyp)} \\
		\midrule
		Komparator echa & AC0, próg z~DAC0 & \textbf{AC1, próg z~DAC1} (PB4=AINP3 jest na AC1, nie AC0 --- Tab.~5-1) \\
		Generacja TX & TCA0 WO5 (sprzęt.\ PWM, PA5) & Programowe taktowanie burstu na \texttt{PB0} \\
		Pin LED & PC0 & \textbf{PA4} \\
		Pin przycisku & PA7 & \textbf{PA3} \\
		Pin NTC & PB1 (AIN10) & \textbf{PA5} (AIN5) \\
		Pin echa & PB0 (AIN11)/PB4 & \textbf{PB4} (przeprowadzony z~PA2) \\
		USART TXD & PA1, \SI{9600}{\baud} & \textbf{PB2, \SI{115200}{\baud}} \\
		Pomiar fazowy (Goertzel) & planowany & \textbf{zaimplementowany} (sek.~\ref{sec:goertzel}) \\
		Detekcja ToF & AC0 capture & ADC obwiednia (czulsza) + AC1 zweryfikowane \\
		Uśpienie & Standby+RTC (plan) & \textbf{zaimplementowane} (Standby+RTC/PIT) \\
		Bramkowanie op-amp (Q1) & Scenariusz 3 (plan) & \textbf{zaimplementowane} \\
		\bottomrule
	\end{tabularx}
\end{table}

\subsection{Uzasadnienie kluczowej korekty (AC0 $\to$ AC1)}

Tabela~5-1 noty katalogowej \textit{ATtiny1614/1616/1617} (DS40002204A, str.~18)
jednoznacznie wskazuje, że pin \texttt{PB4} pełni funkcję \texttt{AINN1}
(wejście \emph{ujemne}) dla komparatora \texttt{AC0}, natomiast funkcję
\texttt{AINP3} (wejście \emph{dodatnie}) dopiero dla komparatora \texttt{AC1}.
Ponieważ sygnał echa wymaga wejścia dodatniego, sprzętowa detekcja musi
korzystać z~\texttt{AC1}. Zgodnie z~architekturą układu (datasheet sek.~29:
,,instance~$n$ of the AC uses instance~$n$ of the DAC'') próg dla \texttt{AC1}
pochodzi z~\texttt{DAC1}, którego referencję ustawiają pola \texttt{DAC1REFSEL}
(\texttt{VREF.CTRLC}) oraz \texttt{DAC1REFEN} (\texttt{VREF.CTRLB}). Korekta ta
była warunkiem uzyskania niezerowych zatrzaśnięć licznika \texttt{TCB0}.

% ══════════════════════════════════════════════════════════════
\section{Dostarczane pliki (deliverables)}
\label{sec:deliverables}

Zgodnie z~wymaganiami Etapu~2 paczka zawiera:

\begin{table}[H]
	\centering
	\caption{Zawartość paczki dostarczanej Komisji}
	\label{tab:deliverables}
	\begin{tabularx}{\textwidth}{lX}
		\toprule
		\textbf{Katalog / plik} & \textbf{Zawartość} \\
		\midrule
		\texttt{firmware/main.c} & Źródło oprogramowania (AVR-GCC, rejestrowo) \\
		\texttt{firmware/main.hex} & Plik wynikowy (Intel HEX) do wgrania przez UPDI \\
		\texttt{firmware/Makefile} & Skrypt budowania i~programowania \\
		\texttt{cad/} & Źródła CAD: \texttt{bezkontaktowy\_miernik\_odległości.sch} + \texttt{.brd} (EAGLE) \\
		\texttt{gerber/} & Komplet Gerber (zip~v17): copper/mask/silk/paste + profil + drill (\texttt{.xln}) + PnP \\
		\texttt{dokumentacja-etap2/} & Niniejsza dokumentacja (\texttt{.tex} + \texttt{.pdf}) \\
		\texttt{*.png} & Rendery PCB i~obudowy, schemat \\
		\bottomrule
	\end{tabularx}
\end{table}

Wszystkie wymagane pliki (prototyp, Gerber, źródła CAD, oprogramowanie
źródłowe i~wynikowe oraz dokumentacja) są kompletne i~dołączone do paczki.
Komplet Gerber wygenerowano zadaniem CAM dla warstw
\texttt{TOP/BOTTOM/soldermask/silkscreen/solderpaste/profile/drill}.

% ══════════════════════════════════════════════════════════════
\section{Dalszy rozwój}
\label{sec:dalsze}

Kierunki dalszego rozwoju prototypu:

\begin{itemize}[nosep]
	\item \textbf{Pełne wykorzystanie fazy do sub-mm} --- etap fazowy (Goertzel)
	jest zaimplementowany i~mierzalnie poprawny ($\sim83^\circ/\si{\milli\metre}$),
	jednak pewny wybór bina $\lambda/2$ wymaga zmniejszenia jittera pomiaru
	zgrubnego poniżej $\lambda/4\approx\SI{2,2}{\milli\metre}$ (obecnie
	${\sim}\SI{4}{\milli\metre}$). Kierunki: silniejszy nadajnik (push-pull),
	uśrednianie fazy z~wielu pingów, interpolacja onsetu.
	\item \textbf{Kalibracja wielopunktowa w~EEPROM} --- obecnie kalibracja
	liniowa (dwa współczynniki w~firmware).
	\item \textbf{Sterowanie push-pull nadajnika} (\SI{+6}{\deci\bel} SNR) ---
	wymaga modyfikacji schematu (usunięcie BC817); poprawia SNR ech dalekich.
\end{itemize}

% ══════════════════════════════════════════════════════════════
\section{Podsumowanie}
\label{sec:podsumowanie}

Zbudowany prototyp realizuje bezkontaktowy pomiar odległości w~zakresie
\SIrange{10}{50}{\centi\metre} metodą hybrydową Pulse-Phase Ranging:
pomiar zgrubny Time-of-Flight przez próbkowanie obwiedni echa (ADC) oraz
pomiar fazy nośnej \SI{40}{\kilo\hertz} algorytmem Goertzela, z~kompensacją
temperatury NTC. Dodatkowo zweryfikowano sprzętową ścieżkę detekcji
\texttt{AC1+DAC1+EVSYS+TCB0}. Zaimplementowano agresywne zarządzanie mocą
(Standby + RTC, bramkowanie toru analogowego Q1), co zapewnia średni pobór
mocy poniżej \SI{1}{\milli\watt}. Pomiar jest stabilny i~powtarzalny
($\sigma\approx\SI{4}{\milli\metre}$), a~faza zmienia się zgodnie z~teorią
($\sim83^\circ/\si{\milli\metre}$). Wszystkie wymagania konkursowe (napięcie,
zakres, obiekt, koszt) są spełnione.

\end{document}
